\section{Các tác nhân}
Hệ thống có 4 tác nhân chính là: Khách, Người nội trợ, Thành viên gia đình và Quản trị viên (Admin). Khách là vai trò của người dùng khi chưa có tài khoản hoặc chưa đăng nhập vào hệ thống. Người nội trợ là vai trò của người dùng chính, trực tiếp tạo tài khoản, thiết lập dữ liệu và quản lý các hoạt động mua sắm, ăn uống của gia đình. Thành viên gia đình là những người dùng được Người nội trợ mời vào nhóm để cùng phối hợp sử dụng thực phẩm và đi chợ. Quản trị viên là người vận hành, kiểm soát dữ liệu nền tảng và quản lý toàn bộ người dùng trong hệ thống.

\section{Biểu đồ use case tổng quan}
Khi chưa đăng nhập, Khách có thể thực hiện đăng ký tài khoản mới và đăng nhập vào hệ thống. Sau khi đăng nhập thành công: Người nội trợ có toàn quyền quản lý đối với không gian lưu trữ của gia đình. Họ có thể: quản lý danh sách mua sắm (thêm, đánh dấu đã mua), phân loại thực phẩm, quản lý thực phẩm tồn kho (trong tủ lạnh), nhận cảnh báo khi sắp hết hạn sử dụng, lên kế hoạch bữa ăn, xem báo cáo lãng phí/tiêu thụ và quản lý nhóm gia đình (tạo nhóm, mời thành viên, phân quyền). Thành viên gia đình có thể tương tác với kho thực phẩm chung bằng cách cập nhật số lượng sử dụng thực phẩm (trừ kho) và chia sẻ, đồng bộ danh sách mua sắm với Người nội trợ để tránh mua trùng lặp.

Đối với Quản trị viên (Admin), họ có các chức năng quản lý cấp cao bao gồm: quản lý danh mục (loại thực phẩm, công thức mẫu), kiểm soát nội dung và người dùng (duyệt công thức, khóa/mở khóa tài khoản), quản lý danh sách tài khoản người dùng và theo dõi hiệu suất hoạt động của toàn hệ thống. Các use case của quản trị viên trong biểu đồ use case tổng quan này là use case phức hợp của một nhóm các use case. Chi tiết về các use case phức này được đưa ra trong các biểu đồ phân rã ở phần sau.

\section{Quy trình nghiệp vụ}
Trong hệ thống Đi chợ tiện lợi, có 4 quy trình nghiệp vụ chính tương ứng với các nhóm chức năng cốt lõi của Người nội trợ và Thành viên gia đình: Quy trình quản lý danh sách mua sắm, Quy trình quản lý thực phẩm trong tủ lạnh, Quy trình gợi ý món ăn và Quy trình lên kế hoạch bữa ăn. Bốn quy trình này có mối liên hệ chặt chẽ với nhau: danh sách mua sắm sau khi hoàn tất sẽ tự động cập nhật vào kho tủ lạnh; dữ liệu trong tủ lạnh là nguồn để gợi ý món ăn và lên kế hoạch bữa ăn; còn kế hoạch bữa ăn khi thiếu nguyên liệu sẽ sinh ra danh sách mua sắm mới, tạo thành một vòng khép kín hỗ trợ người dùng quản lý hiệu quả việc đi chợ và tiêu dùng thực phẩm trong gia đình.

\subsection{Quy trình quản lý danh sách mua sắm}
Người nội trợ tạo một danh sách mua sắm mới theo ngày hoặc theo tuần và đặt tên cho danh sách. Tiếp đó, thêm các thực phẩm cần mua kèm số lượng vào danh sách; hệ thống tự động phân loại các thực phẩm này theo nhóm danh mục. Nếu cần, Người nội trợ có thể chỉnh sửa nhóm thủ công. Sau khi hoàn thiện, có thể chia sẻ danh sách với các Thành viên. Trong quá trình mua hàng, cập nhật trạng thái ``Đã mua'' cho từng món; trạng thái này được đồng bộ thời gian thực. Sau khi mua hàng xong, hệ thống tự động chuyển toàn bộ thực phẩm đã mua vào kho tủ lạnh, ghi nhận lịch sử mua sắm.

\subsection{Quy trình quản lý thực phẩm trong tủ lạnh}
Sau khi mua hàng, hệ thống tự động đưa các thực phẩm đã mua vào kho tủ lạnh; ngoài ra, có thể nhập thủ công các thực phẩm khác. Hằng ngày, hệ thống tự động quét kho và gửi thông báo nhắc nhở cho Người nội trợ về những thực phẩm sắp hết hạn (trước 3 ngày), đánh dấu màu nổi bật. Khi chế biến món ăn, cập nhật số lượng thực phẩm đã sử dụng; hệ thống tự động trừ kho theo thời gian thực và cảnh báo nếu nguyên liệu cạn. Người dùng cũng có thể tìm kiếm, chỉnh sửa hoặc xoá thực phẩm trong kho.

\subsection{Quy trình gợi ý món ăn}
Người nội trợ truy cập tính năng gợi ý món ăn từ trang chủ. Hệ thống kiểm tra danh sách thực phẩm trong tủ lạnh, đối chiếu với kho công thức nấu ăn để tìm ra những món có thể chế biến, ưu tiên những món sử dụng nguyên liệu sắp hết hạn. Nếu món ăn yêu cầu nguyên liệu mà tủ lạnh chưa có đủ, hệ thống hiển thị danh sách còn thiếu và cho phép thêm nhanh vào danh sách mua sắm. Sau khi chế biến, xác nhận đã nấu món ăn để hệ thống tự động trừ kho.

\subsection{Quy trình lên kế hoạch bữa ăn}
Người nội trợ chọn chế độ lên kế hoạch theo ngày hoặc theo tuần và mở giao diện lịch thực đơn. Với mỗi bữa ăn, hệ thống đề xuất thực đơn phù hợp dựa trên thực phẩm hiện có và sở thích. Sau khi xác nhận thực đơn, hệ thống đối chiếu nguyên liệu. Nếu thiếu, hệ thống tự động sinh ra một danh sách mua sắm bổ sung. Khi đến ngày thực hiện, đánh dấu hoàn tất bữa ăn để hệ thống cập nhật lịch sử.
